%Travail des forces s'exerçant sur un pendule
Un pendule est constitué par une bille métallique de masse $m=\qty{250}{\g}$ et
de dimensions négligeables attachée à l'extrémité d'un fil de longueur
$\ell=\qty{1.00}{\m}$. L'autre extrémité du fil est fixée en un point fixe $O$.

Le pendule étant dans sa position d'équilibre, on l'écarte, fil tendu, d'un
angle $\theta_0=\ang{60}$ et on le lâche sans vitesse initiale.
\begin{enumerate}
    \item Dresser un bilan des forces s'appliquant sur la bille du pendule en
          mouvement. Représenter ces forces sur un schéma.
    \item Calculer le travail de ces forces entre la position initiale du
          pendule et le premier passage à la position d'équilibre.
\end{enumerate}
\begin{donnees}
    Intensité de la pesanteur~: $g=\qty{9.8}{\N\per\kg}$.
\end{donnees}
